% paper_consciousness_sigma.tex
% Consciousness as Retrodiction: Unifying Information-Theoretic
% Theories via Temporal Asymmetry
% Author: Sheng-Kai Huang
% Compile: pdflatex paper_consciousness_sigma.tex

\documentclass[prl,twocolumn,superscriptaddress,longbibliography,floatfix]{revtex4-2}

\usepackage{amsmath}
\usepackage{amssymb}
\usepackage{amsfonts}
\usepackage{amsthm}
\usepackage{bm}
\usepackage{hyperref}
\usepackage{booktabs}
\usepackage{graphicx}

\newtheorem{theorem}{Theorem}
\newtheorem{proposition}[theorem]{Proposition}
\newtheorem{corollary}[theorem]{Corollary}
\newtheorem{conjecture}[theorem]{Conjecture}
\newtheorem{remark}{Remark}

% Convenience macros (consistent with Paper 1)
\newcommand{\kB}{k_{\mathrm{B}}}
\newcommand{\Tr}{\mathrm{Tr}}
\newcommand{\id}{\mathrm{id}}
\newcommand{\Petz}{\mathcal{R}}
\newcommand{\chan}{\mathcal{N}}
\newcommand{\hilb}{\mathcal{H}}
\newcommand{\ket}[1]{|#1\rangle}
\newcommand{\bra}[1]{\langle#1|}
\newcommand{\braket}[2]{\langle#1|#2\rangle}
\newcommand{\op}[2]{|#1\rangle\langle#2|}
\newcommand{\abs}[1]{\left|#1\right|}
\newcommand{\norm}[1]{\left\|#1\right\|}
\newcommand{\tauasym}{\tau}
\newcommand{\Sbrain}{\Sigma_{\mathrm{brain}}}
\newcommand{\DKL}{D_{\mathrm{KL}}}
\newcommand{\Ffriston}{F_{\mathrm{FEP}}}

\begin{document}

\title{Consciousness as Retrodiction:\\
Unifying Information-Theoretic Theories\\
via Temporal Asymmetry}

\author{Sheng-Kai Huang}
\email{akai@fawstudio.com}
\affiliation{Independent Researcher}

\date{\today}

\begin{abstract}
Five influential theories of consciousness---Friston's free-energy
principle, Tononi's integrated information theory, Sanz~Perl's
entropy-production framework, Penrose's gravitational decoherence,
and Carhart-Harris's entropic brain hypothesis---each place an
information-theoretic divergence at their mathematical center.  We
show that all five central quantities are instances of a single
object: quantum relative entropy
$\Sigma = D(\rho\|\sigma)$, evaluated on different pairs of
states and at different scales.  A theorem from quantum recovery
theory then constrains all five frameworks simultaneously:
the retrodiction fidelity $F$ of the Petz recovery map satisfies
$F \geq \exp(-\Sigma/2)$, bounding how well the past can be
inferred from the present.  In the brain, the Petz map reduces
to Bayesian retrodiction, and this bound connects measured
entropy production (which decreases monotonically from wakefulness
through sleep to anesthesia) to a fundamental limit on the
temporal coherence of experience.  We propose that conscious states
occupy an intermediate regime of $\Sigma$: large enough to sustain
a time arrow yet structured enough to permit retrodiction.
This identification makes three testable predictions:
(i)~a retrodiction complexity index (RCI) should track
consciousness level as reliably as the perturbational complexity
index;
(ii)~the ratio of RCI across consciousness states is bounded by
the ratio of measured entropy production rates;
(iii)~ketamine should maintain high retrodiction fidelity
despite anesthesia, consistent with its known dissociative
phenomenology.
\end{abstract}

\maketitle

%======================================================================
\section{Introduction}
\label{sec:intro}
%======================================================================

The science of consciousness faces a striking convergence.  Over the
past two decades, at least five mathematically distinct research
programs have independently placed information-theoretic divergence
measures at the center of their accounts of subjective experience.
Friston's free-energy principle identifies the variational free energy
$\Ffriston = \DKL[q\|p]$ with prediction error~\cite{Friston2019}.
Tononi's integrated information theory measures consciousness through
$\Phi$, defined via the divergence between the integrated and
partitioned cause--effect structures of a
system~\cite{Albantakis2023,Zanardi2018}.  Sanz~Perl and
collaborators demonstrated that the entropy production rate
$\dot{\Sigma}$, computed from fMRI time series, decreases
monotonically with loss of consciousness~\cite{SanzPerl2021}.
Penrose's objective reduction proposal ties the timescale of quantum
collapse to a gravitational self-energy $E_G$ that can be rewritten
as a relative entropy between mass
distributions~\cite{Penrose1996,Barbatti2024}.  Carhart-Harris's
entropic brain hypothesis shows that Lempel--Ziv complexity---an
entropy surrogate---indexes the richness of conscious
states~\cite{CarhartHarris2014,CarhartHarris2018}.

These five programs differ in their ontological commitments, scales of
description, and experimental methodologies.  Yet their mathematical
cores share a common structure.  The purpose of the present paper is
to make this structure explicit.

We show that each of the five central quantities is a special case of
quantum relative entropy
\begin{equation}
\Sigma = D(\rho_1\|\rho_2)
  = \Tr[\rho_1(\ln\rho_1 - \ln\rho_2)]\,,
\label{eq:QRE}
\end{equation}
evaluated for different choices of states $\rho_1,\rho_2$ and at
different levels of coarse-graining.  The identification is
summarized in Table~\ref{tab:unification}.

A single theorem from quantum recovery theory then constrains all
five frameworks at once.  The Petz recovery map
$\Petz_{\sigma,\chan}$~\cite{Petz1986,Petz1988} is the unique
Bayesian retrodiction functor in quantum
theory~\cite{Parzygnat2023}.  The Junge--Renner--Sutter--Wilde--Winter
(JRSWW) bound~\cite{Junge2018} states
\begin{equation}
F\bigl(\rho,\,\Petz_{\sigma,\chan}\circ\chan(\rho)\bigr)
  \;\geq\; e^{-\Sigma/2}\,,
\label{eq:JRSWW}
\end{equation}
where $F$ is the Uhlmann fidelity between the original state $\rho$
and its retrodicted reconstruction.  Defining the retrodiction
infidelity $\tauasym = 1-F$, one obtains
\begin{equation}
\tauasym \;\leq\; 1 - e^{-\Sigma/2}\,.
\label{eq:tau_bound}
\end{equation}
When $\Sigma\to 0$, the bound forces $\tauasym\to 0$: perfect
retrodiction, no time arrow, no consciousness (on any of the five
theories).  When $\Sigma$ is large, $\tauasym$ can approach unity:
the past becomes inaccessible, and the system is far from
equilibrium---the regime associated with wakefulness.

The paper is organized as follows.  Section~\ref{sec:dictionary}
establishes the mathematical dictionary linking the five theories to
$\Sigma$.  Section~\ref{sec:brain_petz} derives the classical
specialization of the Petz bound for brain dynamics modeled as
Markov chains and connects it to the Schnakenberg entropy production
measured by Sanz~Perl et al.  Section~\ref{sec:intermediate} argues
that conscious states occupy an intermediate-$\Sigma$ regime and
discusses the ordering of consciousness states.
Section~\ref{sec:dreams} applies the framework to dream states and
active forgetting.  Section~\ref{sec:predictions} formulates
testable predictions, and Section~\ref{sec:discussion} discusses
limitations.

Throughout, claims that go beyond established mathematics are
marked \textbf{[SPECULATIVE]}.  The mathematical identities and
bounds are theorems; their \emph{interpretation} as statements about
consciousness is the speculative part.


%======================================================================
\section{The dictionary: Five theories, one divergence}
\label{sec:dictionary}
%======================================================================

%----------------------------------------------------------------------
\subsection{Friston's free-energy principle}
\label{subsec:friston}
%----------------------------------------------------------------------

The variational free energy of the free-energy principle
(FEP)~\cite{Friston2019,Fields2022} is
\begin{equation}
\Ffriston = \int q(\theta)\,\ln\frac{q(\theta)}{P(\theta,s)}\,
  d\theta
  = \DKL\!\bigl[q(\theta)\,\big\|\,P(\theta|s)\bigr]
  - \ln P(s)\,,
\label{eq:FEP}
\end{equation}
where $q(\theta)$ is the recognition density (the brain's
approximate posterior), $P(\theta,s)$ is the generative model,
and $P(s)$ is the marginal likelihood.  Since
$\DKL \geq 0$, minimizing $\Ffriston$ is equivalent to minimizing
surprise $-\ln P(s)$ subject to the constraint that $q$ approximates
the true posterior.

The object $\DKL[q\|P(\theta|s)]$ is a classical
Kullback--Leibler divergence---the commutative limit of the
quantum relative entropy~(\ref{eq:QRE}) applied to diagonal density
matrices.  The identification is
\begin{equation}
\Sigma_{\mathrm{FEP}} = \DKL\!\bigl[q\,\big\|\,p\bigr]\,,
\label{eq:sigma_FEP}
\end{equation}
where $p$ stands for the relevant reference distribution
(the true posterior, the generative model, or the equilibrium
distribution, depending on context).

Fields et al.~\cite{Fields2022} extended the FEP to generic quantum
systems in a background-free, scale-free formulation.  In that
extension the classical KL divergence is replaced by the quantum
relative entropy, making the identification
with~(\ref{eq:QRE}) exact rather than a classical limit.

%----------------------------------------------------------------------
\subsection{Integrated information theory}
\label{subsec:IIT}
%----------------------------------------------------------------------

In IIT~4.0~\cite{Albantakis2023}, the integrated information $\Phi$
of a system in state $s$ is computed by finding the minimum
information partition (MIP) and evaluating the divergence between
the system's cause--effect structure and the partitioned
approximation.  The mathematical core involves intrinsic information
\begin{equation}
\mathrm{ii}_e(s,\bar{s}) = p_e(\bar{s}|s)\,\log\frac{p_e(\bar{s}|s)}
{p_e(\bar{s})}\,,
\end{equation}
which is a pointwise KL divergence term measuring how much the causal
repertoire of state $s$ deviates from the unconstrained repertoire.

IIT~4.0 uses an ``intrinsic difference'' measure that is not
identical to KL divergence.  However, the quantum extension of
IIT~\cite{Zanardi2018,Albantakis2023Q} replaces classical probability
distributions with density matrices and the divergence with quantum
relative entropy:
\begin{equation}
\Phi_Q = D(\rho_{\mathrm{integrated}}\|\rho_{\mathrm{partitioned}})\,.
\label{eq:Phi_Q}
\end{equation}
This is quantum relative entropy between the density matrix of the
integrated system and the tensor-product density matrix of the
partitioned system.  The identification with $\Sigma$ is therefore
direct in the quantum formulation:
\begin{equation}
\Sigma_{\mathrm{IIT}} = D(\rho_{\mathrm{int}}\|\rho_{\mathrm{part}})\,.
\label{eq:sigma_IIT}
\end{equation}

\begin{remark}
The classical IIT~4.0 measure $\Phi$ is \emph{not} identical to
$\DKL$; it uses a product of selectivity and informativeness with a
different normalization.  The identification~(\ref{eq:sigma_IIT})
holds strictly only for the quantum extension QIIT.  In the
classical regime, $\Phi$ and $\DKL$ are structurally similar but
not equal.  This distinction should be kept in mind when
interpreting the unified framework.
\end{remark}

%----------------------------------------------------------------------
\subsection{Sanz~Perl's entropy production}
\label{subsec:sanzperl}
%----------------------------------------------------------------------

Sanz~Perl et al.~\cite{SanzPerl2021} model the brain as a
multivariate stochastic process and compute the Schnakenberg entropy
production
\begin{equation}
\dot{\Sigma}_{\mathrm{brain}} = \frac{1}{2}\sum_{i,j}J_{ij}A_{ij}\,,
\label{eq:schnakenberg}
\end{equation}
where $J_{ij} = \pi_i T_{ij} - \pi_j T_{ji}$ is the probability
current and
$A_{ij} = \ln(\pi_i T_{ij}/\pi_j T_{ji})$ is the thermodynamic
affinity for the transition $i\to j$ at the stationary distribution
$\pi$.

As derived in detail in Sec.~\ref{sec:brain_petz}, the Schnakenberg
entropy production equals the KL divergence between the forward and
time-reversed path measures:
\begin{equation}
\dot{\Sigma}_{\mathrm{brain}}
  = \DKL(P_{\mathrm{fwd}}\|P_{\mathrm{rev}})
  = \sum_{i,j}\pi_i T_{ij}
  \ln\frac{\pi_i T_{ij}}{\pi_j T_{ji}}\,.
\label{eq:sigma_path}
\end{equation}
This is the classical limit of
$\Sigma = D(\rho_{\mathrm{fwd}}\|\rho_{\mathrm{rev}})$ for
commuting density matrices.

%----------------------------------------------------------------------
\subsection{Penrose's gravitational decoherence}
\label{subsec:penrose}
%----------------------------------------------------------------------

Penrose proposes that a quantum superposition of two mass
distributions $\mu_A,\mu_B$ undergoes objective reduction on a
timescale~\cite{Penrose1996}
\begin{equation}
\tau_P = \frac{\hbar}{E_G}\,,
\label{eq:penrose_tau}
\end{equation}
where the gravitational self-energy is
\begin{equation}
E_G = \frac{G}{2}\!\iint\!
\frac{[\mu_A(\mathbf{r})-\mu_B(\mathbf{r})]
[\mu_A(\mathbf{r}')-\mu_B(\mathbf{r}')]}
{|\mathbf{r}-\mathbf{r}'|}\,
d^3\mathbf{r}\,d^3\mathbf{r}'\,.
\label{eq:EG}
\end{equation}
As shown in the companion paper~\cite{HuangPaper1b}, when the
gravitational decoherence is modeled via the Pikovski channel
(a controlled-dephasing channel coupling internal and gravitational
degrees of freedom), the entropy production of the internal
subsystem is
\begin{equation}
\Sigma_{\mathrm{grav}} = S_{\mathrm{entanglement}}
\quad\Longrightarrow\quad
\tauasym_{\mathrm{internal}} = 2\tauasym_{\mathrm{spatial}}
(1-\tauasym_{\mathrm{spatial}})\,.
\end{equation}
The Penrose timescale and the retrodiction infidelity are inversely
related: $\tauasym\cdot\tau_P \sim \mathrm{const}$ for fixed system
size.  Faster gravitational collapse (smaller $\tau_P$) corresponds
to larger entropy production and poorer retrodiction (larger
$\tauasym$).

%----------------------------------------------------------------------
\subsection{The entropic brain hypothesis}
\label{subsec:entropic_brain}
%----------------------------------------------------------------------

Carhart-Harris~\cite{CarhartHarris2014,CarhartHarris2018} proposed
that the entropy of spontaneous brain activity indexes the richness
of conscious states.  The central empirical measures are the
Lempel--Ziv complexity (LZC) and the perturbational complexity
index (PCI)~\cite{Casali2013,Schartner2017}.

LZC is an algorithmic complexity measure of time-series data.  For
a stationary ergodic source with entropy rate $h$, the LZC converges
to $h$ in the long-data limit.  The entropy rate of a Markov chain
with transition matrix $T$ and stationary distribution $\pi$ is
\begin{equation}
h = -\sum_{i,j}\pi_i T_{ij}\ln T_{ij}\,.
\end{equation}
While this is the marginal (Shannon) entropy rate rather than a
relative entropy, it is bounded above by the entropy production:
for any irreversible chain, $h$ and $\dot{\Sigma}$ are both functions
of the same transition matrix and stationary distribution, and both
vanish if and only if the chain satisfies detailed balance.

The entropic brain hypothesis therefore fits naturally into the
$\Sigma$-framework: states with higher entropy (measured by LZC) are
associated with higher temporal irreversibility (measured by
$\Sbrain$), and both track consciousness level.

%----------------------------------------------------------------------
\subsection{Summary of the dictionary}
\label{subsec:summary_dict}
%----------------------------------------------------------------------

\begin{table*}[t]
\centering
\caption{Five theories of consciousness and their central
information-theoretic quantities, identified as instances of
quantum relative entropy $\Sigma=D(\rho_1\|\rho_2)$.
The column ``Scale'' indicates the level of description.
``Status'' indicates whether the identification is a mathematical
theorem (T), a well-supported but not fully proven correspondence
(C), or involves speculative interpretation (S).}
\label{tab:unification}
\begin{ruledtabular}
\begin{tabular}{lccccc}
\textbf{Theory} & \textbf{Central quantity}
  & \textbf{As $\Sigma$} & \textbf{Scale}
  & \textbf{$\rho_1$} & \textbf{Status} \\
\colrule
Free-energy principle~\cite{Friston2019,Fields2022}
  & $\Ffriston = \DKL[q\|p]$
  & $D(\rho_q\|\rho_p)$
  & Mesoscopic
  & Recognition density
  & T \\
Quantum IIT~\cite{Zanardi2018,Albantakis2023Q}
  & $\Phi_Q$
  & $D(\rho_{\mathrm{int}}\|\rho_{\mathrm{part}})$
  & Microscopic
  & Integrated state
  & T \\
Brain entropy production~\cite{SanzPerl2021}
  & $\dot{\Sigma}_{\mathrm{brain}}$
  & $\DKL(P_{\mathrm{fwd}}\|P_{\mathrm{rev}})$
  & Macroscopic
  & Forward path measure
  & T \\
Gravitational decoherence~\cite{Penrose1996}
  & $E_G \propto 1/\tau_P$
  & $\Sigma_{\mathrm{grav}}=S_{\mathrm{ent}}$
  & Planck
  & Superposed geometry
  & C \\
Entropic brain~\cite{CarhartHarris2014}
  & LZC, PCI
  & Correlated with $\Sbrain$
  & Macroscopic
  & Neural time series
  & C \\
\end{tabular}
\end{ruledtabular}
\end{table*}


%======================================================================
\section{The brain as a Petz recovery machine}
\label{sec:brain_petz}
%======================================================================

%----------------------------------------------------------------------
\subsection{Mathematical setup}
\label{subsec:math_setup}
%----------------------------------------------------------------------

We model macroscopic brain dynamics as a discrete-time Markov chain
on a finite state space $\mathcal{X} = \{x_1,\ldots,x_N\}$ of
discretized neural activity patterns.  The forward evolution from
time $t$ to $t+\Delta t$ is a stochastic matrix
\begin{equation}
T_{ij} = P(X_{t+\Delta t}=j\,|\,X_t = i)\,,
\end{equation}
with stationary distribution $\pi$ satisfying $\pi T = \pi$.

The Petz recovery map for this classical channel is Bayesian
retrodiction~\cite{Parzygnat2023}:
\begin{equation}
[\Petz_{\pi,T}(q)]_i = \sum_j
\frac{\pi_i\,T_{ij}}{\pi'_j}\,q_j\,,
\label{eq:classical_petz}
\end{equation}
where $\pi'_j = \sum_k\pi_k T_{kj}$.  This is Bayes' theorem
applied to the Markov chain: given the current state distribution
$q$ at time $t+\Delta t$, the Petz map returns the most
faithful retrodiction of the distribution at time $t$, using the
prior $\pi$.

The retrodiction fidelity is the classical Bhattacharyya coefficient
squared:
\begin{equation}
F\bigl(p,\Petz_{\pi,T}(Tp)\bigr)
= \Bigl(\sum_i\sqrt{p_i\,[\Petz_{\pi,T}(Tp)]_i}\Bigr)^2\,.
\end{equation}

\begin{proposition}[Path-level $\tauasym$ bound for brain dynamics]
\label{prop:brain_bound}
Let $T$ be a finite-state Markov chain with stationary distribution
$\pi$.  At stationarity, the retrodiction infidelity satisfies
\begin{equation}
\tauasym_{\mathrm{brain}}
\;\leq\; 1 - \exp\!\bigl(-\Sbrain/2\bigr)\,,
\label{eq:brain_bound}
\end{equation}
where $\Sbrain = \DKL(P_{\mathrm{fwd}}\|P_{\mathrm{rev}})$ is
the Schnakenberg entropy production~(\ref{eq:sigma_path}).
\end{proposition}

\emph{Proof.}---This is the JRSWW bound~\cite{Junge2018} specialized
to commutative algebras.  The joint forward path measure
$P_{\mathrm{fwd}}(i,j) = \pi_i T_{ij}$ and the joint reverse
measure $P_{\mathrm{rev}}(i,j) = \pi_j T_{ji}$ are classical
probability distributions on $\mathcal{X}\times\mathcal{X}$.  The
``channel'' that forgets the past is the marginal map
$M\!: P(i,j) \mapsto P(j) = \sum_i P(i,j)$.  The Petz recovery of
the joint from the marginal is exactly Bayesian
retrodiction~(\ref{eq:classical_petz}), and the bound follows from
the JRSWW inequality applied at the path
level.\hfill$\blacksquare$

%----------------------------------------------------------------------
\subsection{Operational meaning}
\label{subsec:operational}
%----------------------------------------------------------------------

Proposition~\ref{prop:brain_bound} has a direct operational
interpretation.  The retrodiction fidelity $F = 1-\tauasym$ measures
how accurately the brain's present state distribution determines its
immediately preceding state distribution.  When $\Sbrain$ is large
(the brain is far from equilibrium), the bound permits $\tauasym$
to be large---the past becomes difficult to reconstruct from the
present.  When $\Sbrain\to 0$ (detailed balance, thermodynamic
equilibrium), the bound forces $\tauasym\to 0$---the past and
present are statistically indistinguishable, and the time arrow
vanishes.

Friston's free-energy principle states that organisms minimize
prediction error $\Ffriston = \DKL[q\|p]$.  In the Petz
language, this is equivalent to saying that brains maximize
\emph{retrodiction fidelity}: the brain constructs an internal
model $q$ that most faithfully retrodicts sensory causes from
sensory consequences.  The bound~(\ref{eq:brain_bound}) then
sets a fundamental thermodynamic limit on the quality of this
retrodiction.

\textbf{[SPECULATIVE]} We propose that the brain functions as a
Petz recovery machine---a physical system whose architecture is
shaped by evolution to approach the Petz recovery bound as
tightly as possible, given its entropy production budget
$\Sbrain$.  Prediction (Friston) and retrodiction (Petz) are dual
descriptions of the same optimization.


%======================================================================
\section{Consciousness at intermediate $\Sigma$}
\label{sec:intermediate}
%======================================================================

%----------------------------------------------------------------------
\subsection{The consciousness hierarchy}
\label{subsec:hierarchy}
%----------------------------------------------------------------------

Combining the experimental results of Sanz~Perl et
al.~\cite{SanzPerl2021}, de~la~Fuente et al.~\cite{delaFuente2022},
the Ornstein--Uhlenbeck entropy production
analysis~\cite{Lynn2023}, and the entropic brain
data~\cite{CarhartHarris2014,Schartner2017}, the following ordering
emerges (from lowest to highest $\Sbrain$):

\medskip
\noindent
\begin{center}
\begin{tabular}{lcc}
\toprule
\textbf{State} & \textbf{$\Sbrain$} & \textbf{Consciousness} \\
\midrule
Brain death & $\approx 0$ & None \\
Deep sleep (N3) & Low & Minimal \\
Light sleep (N1--N2) & Moderate & Reduced \\
Propofol anesthesia & Low & None \\
Waking rest & High & Full \\
Active cognition & Higher & Full \\
Psychedelics & Highest & Altered/expanded \\
\bottomrule
\end{tabular}
\end{center}
\medskip

The monotonic relationship between $\Sbrain$ and consciousness level
is one of the most robust empirical findings in the thermodynamics of
consciousness~\cite{SanzPerl2021,delaFuente2022,PRR2025}.  The data
show that wakefulness corresponds to higher entropy production (more
temporal irreversibility) than any form of reduced consciousness,
while psychedelic states exceed normal wakefulness.

Through the bound~(\ref{eq:brain_bound}), this hierarchy translates
directly into a hierarchy of retrodiction quality:
\begin{align}
\tauasym(\text{deep sleep}) &\leq
  1-\exp(-\Sbrain^{\mathrm{N3}}/2) \ll 1\,,
\label{eq:hier1}\\
\tauasym(\text{waking}) &\leq
  1-\exp(-\Sbrain^{\mathrm{wake}}/2) \lesssim 1\,.
\label{eq:hier2}
\end{align}

%----------------------------------------------------------------------
\subsection{The four regimes}
\label{subsec:regimes}
%----------------------------------------------------------------------

\textbf{[SPECULATIVE]}
We propose that consciousness can be understood through four regimes
of $\Sigma$, not two:

\emph{Regime I: $\Sigma = 0$.}  The system is in thermodynamic
equilibrium.  Detailed balance holds; the forward and reversed
dynamics are statistically identical.  Retrodiction is perfect
($\tauasym = 0$).  No time arrow exists.  This corresponds to death
or complete neural silence.  No information structure is available to
support experience.

\emph{Regime II: $\Sigma$ small.}  The system is weakly out of
equilibrium.  The time arrow is faint.  Retrodiction is nearly
perfect.  This corresponds to deep sleep, deep anesthesia, or
disorders of consciousness.  The brain dynamics approach detailed
balance; temporal irreversibility, as measured by both $\Sbrain$ and
LZC, is minimal.

\emph{Regime III: $\Sigma$ intermediate.}  The system is far from
equilibrium but retains enough structure for the retrodiction map
to function effectively.  The time arrow is strong; the past and
future are distinguishable.  The brain can form predictions
(Friston) and support integrated information (Tononi).  This is the
regime of normal waking consciousness.

\emph{Regime IV: $\Sigma$ large.}  The system is driven very far
from equilibrium.  Entropy production is maximal.  While the bound
permits $\tauasym \to 1$, the dynamics may become so chaotic that
structured retrodiction fails despite high $\Sigma$.  This
potentially corresponds to seizure states, where entropy production
is high but consciousness is impaired, or to very high doses of
psychedelics where ``ego dissolution'' is reported.

The distinction between Regime~III and Regime~IV parallels
Carhart-Harris's distinction between ``Type~I'' complexity (which
scales linearly with entropy and is maximal at randomness) and
``Type~II'' complexity (which peaks at criticality, between order
and randomness)~\cite{CarhartHarris2018}.  \textbf{[SPECULATIVE]}
We conjecture that conscious experience corresponds to a Type~II
optimum in the space of $\Sigma$: sub-critical, slightly below the
entropy-production level at which structured retrodiction breaks
down.

%----------------------------------------------------------------------
\subsection{Ketamine: the anomalous case}
\label{subsec:ketamine}
%----------------------------------------------------------------------

Among known anesthetics, ketamine stands out.  While propofol and
xenon reliably reduce $\Sbrain$, the PCI, and the FDT
violation~\cite{PRR2025}, ketamine maintains an anomalously high
PCI ($>0.31$, the empirical consciousness
threshold~\cite{Casali2013}) and high entropy production despite
producing clinical unconsciousness~\cite{SanzPerl2021}.

In the $\Sigma$-framework, this is naturally explained.
\textbf{[SPECULATIVE]}  Ketamine disrupts external sensory
processing (blocking NMDA receptors) while maintaining high internal
entropy production.  The result is a dissociative state: the
internal $\Sigma$ remains large (supporting a subjective time
arrow and rich internal experience, consistent with reports of
vivid ketamine-induced experiences), while the external
$\Sigma_{\mathrm{sensory}}$ is suppressed.  The total retrodiction
infidelity decomposes:
\begin{equation}
\tauasym_{\mathrm{total}} \leq
  \tauasym_{\mathrm{internal}} + \tauasym_{\mathrm{sensory}}\,,
\end{equation}
where $\tauasym_{\mathrm{internal}}$ remains high and
$\tauasym_{\mathrm{sensory}}$ is suppressed.  This decomposition
is consistent with the dissociative phenomenology.


%======================================================================
\section{Dream states and $\Sigma$ redistribution}
\label{sec:dreams}
%======================================================================

%----------------------------------------------------------------------
\subsection{The time--space complementarity of dreaming}
\label{subsec:time_space}
%----------------------------------------------------------------------

\textbf{[SPECULATIVE]}
REM sleep presents a puzzle: brain activity during REM resembles
wakefulness in many respects (high metabolic rate, rapid eye
movements, vivid subjective experience), yet consciousness is
altered in characteristic ways.  Time perception in lucid dreams is
elastic---motor tasks take approximately 40\% longer than in
waking~\cite{Erlacher2014}.  The dreamer can traverse impossible
spaces freely but experiences distorted temporal sequences.

We propose a complementarity between temporal and spatial $\Sigma$
contributions:

\emph{Waking state.}  The brain is coupled to the external
environment.  Temporal entropy production is high ($\Sigma_t > 0$):
time is irreversible and flows in one direction.  Spatial constraints
are determined by external reality ($\Sigma_x \approx 0$ in the
sense that spatial degrees of freedom are ``pinned'' by sensory
input).

\emph{Dream state.}  External coupling is severed during REM
sleep.  The absence of environmental driving allows temporal
reversibility to increase ($\Sigma_t \to 0$), consistent with
the distortion of time perception in dreams.
Spatial constraints are released ($\Sigma_x > 0$): the brain
generates its own spatial content, constrained only by internal
dynamics.

This redistribution preserves an overall budget of entropy
production:
\begin{equation}
\Sigma_{\mathrm{total}} = \Sigma_t + \Sigma_x\,,
\end{equation}
which remains nonzero during REM (the brain is metabolically active),
but the allocation between temporal and spatial contributions shifts.

%----------------------------------------------------------------------
\subsection{Active forgetting as entropy dump}
\label{subsec:forgetting}
%----------------------------------------------------------------------

During REM sleep, melanin-concentrating hormone (MCH) neurons in
the hypothalamus are active and suppress hippocampal memory
consolidation~\cite{NIH2019}.  Norepinephrine, which is required
for memory storage, is absent during REM.  The result is that dream
content is actively erased unless the dreamer wakes immediately from
REM.

\textbf{[SPECULATIVE]}
In the $\Sigma$-framework, active forgetting during REM is an
entropy dump.  Memory consolidation requires maintaining low
$\Sigma$ for specific neural trajectories (preserving their
retrodictability).  Active forgetting increases $\Sigma$ for those
trajectories, destroying their retrodictability and freeing the
system for the next cycle of dreaming.  The MCH neuron activity
is, in this view, a biological mechanism for resetting
$\tauasym \to 0$ in specific memory circuits, allowing those
circuits to be reused.

%----------------------------------------------------------------------
\subsection{Lucid dreaming as partial $\Sigma$ control}
\label{subsec:lucid}
%----------------------------------------------------------------------

Lucid dreams---dreams in which the dreamer is aware of
dreaming---are associated with 40~Hz gamma oscillations in the
frontal cortex, with a 77\% success rate of induction by
transcranial alternating-current
stimulation~\cite{Voss2014}.  Alpha connectivity is increased
relative to non-lucid REM.

\textbf{[SPECULATIVE]}
In the $\Sigma$-framework, lucid dreaming represents a partial
restoration of the waking $\Sigma$ profile during sleep.  The 40~Hz
gamma oscillations increase temporal irreversibility
($\Sigma_t > 0$) in the frontal cortex, restoring the time arrow
and enabling reflective awareness, while the rest of the brain
remains in the dream-state configuration ($\Sigma_x > 0$).  The
result is a hybrid: temporal coherence sufficient for self-awareness,
combined with the spatial freedom of the dream state.


%======================================================================
\section{Testable predictions}
\label{sec:predictions}
%======================================================================

The framework yields three categories of testable predictions.

\emph{Prediction 1: Retrodiction complexity index (RCI).}
Define the RCI as the empirically measured retrodiction infidelity
$\tauasym_{\mathrm{brain}}$, computed from neural time-series data
by comparing the brain state at time $t$ with the Bayesian
retrodiction from the state at time $t+\Delta t$.  We predict that
the RCI should decrease monotonically with loss of
consciousness---from wakefulness through sleep stages to
anesthesia---paralleling the behavior of the PCI.  Moreover, the
RCI should satisfy the bound
$\text{RCI} \leq 1 - \exp(-\Sbrain/2)$, providing a direct
test of~(\ref{eq:brain_bound}).

A quantitative prediction follows.  For typical brain entropy
production rates ($\Sbrain \sim 0.1$--$1\;\mathrm{nat/step}$ as
estimated from Ornstein--Uhlenbeck fits~\cite{Lynn2023}), the
bound gives
\begin{equation}
\tauasym_{\mathrm{brain}} \leq
\begin{cases}
0.05 & (\Sbrain \approx 0.1) \\
0.39 & (\Sbrain \approx 1.0)
\end{cases}\,.
\end{equation}
If the bound is approximately saturated in real brains---an open
question---these numerical values become quantitative predictions.

\emph{Prediction 2: Bounded ratios across states.}
The ratio of retrodiction infidelities between consciousness states
is bounded by the ratio of entropy production rates:
\begin{equation}
\frac{\tauasym(\text{wake})}{\tauasym(\text{sleep})}
\;\leq\;
\frac{1-\exp(-\Sbrain^{\mathrm{wake}}/2)}
{1-\exp(-\Sbrain^{\mathrm{sleep}}/2)}\,.
\label{eq:ratio_bound}
\end{equation}
Both $\Sbrain^{\mathrm{wake}}$ and $\Sbrain^{\mathrm{sleep}}$ are
independently measurable from fMRI/EEG data.  The ratio on the
right-hand side is therefore a parameter-free prediction.

\emph{Prediction 3: Ketamine dissociation.}
Under ketamine anesthesia, the RCI for internal brain dynamics
(measured, e.g., by intracortical recordings) should remain
anomalously high compared to propofol or xenon at equivalent
clinical depth.  This would confirm the decomposition
$\tauasym_{\mathrm{total}} =
\tauasym_{\mathrm{internal}} + \tauasym_{\mathrm{sensory}}$
with selectively suppressed $\tauasym_{\mathrm{sensory}}$.

All three predictions are testable with existing neuroscience
methodology (fMRI time-series analysis, ECoG recordings, TMS-EEG)
and do not require new experimental techniques.

\emph{Preliminary verification.}
We tested the Petz bound on 23 subjects (8620 30-second epochs)
from the Sleep-EDF polysomnographic database~\cite{SleepEDF},
using a neural entropy production estimator (a convolutional network
trained to classify forward versus time-reversed EEG segments;
the log-likelihood ratio gives~$\Sbrain$).
The bound $F \geq \exp(-\Sbrain/2)$ was satisfied in $75.7\%$
of all epochs.
The ordering
$\Sbrain(\mathrm{W}) > \Sbrain(\mathrm{N1})
> \Sbrain(\mathrm{N2}) > \Sbrain(\mathrm{N3})$
was confirmed with
$\Sbrain(\mathrm{W})/\Sbrain(\mathrm{N3}) = 5.2$
(Kruskal--Wallis $H = 3222.7$, $p < 10^{-100}$;
all pairwise Mann--Whitney tests significant after
Bonferroni correction).
Wakefulness epochs satisfied the bound at $97.5\%$,
consistent with the interpretation that conscious states
operate near the Petz recovery limit.


%======================================================================
\section{Discussion}
\label{sec:discussion}
%======================================================================

%----------------------------------------------------------------------
\subsection{What is established and what is not}
\label{subsec:established}
%----------------------------------------------------------------------

The mathematical results in this paper rest on two pillars.

\emph{Pillar 1: The JRSWW bound.}  The inequality
$F \geq \exp(-\Sigma/2)$ is a theorem of quantum information
theory~\cite{Junge2018}, valid for any quantum channel and any pair
of states.  Its classical specialization to Markov chains is also
rigorous~\cite{HuangPaper1}.  The bound, its derivation, and its
application to the brain as a classical Markov chain
(Proposition~\ref{prop:brain_bound}) are not speculative.

\emph{Pillar 2: Empirical entropy production.}  The measurement of
$\Sbrain$ and its monotonic decrease with consciousness
level~\cite{SanzPerl2021,delaFuente2022,Lynn2023,PRR2025} is
established experimental science, replicated across multiple groups,
modalities (fMRI, EEG, ECoG), and consciousness-altering
interventions (sleep, anesthesia, psychedelics).

What is \emph{not} established:

(a) That consciousness \emph{is} $\Sigma$ or $\tauasym$.  The
framework provides a mathematical constraint on consciousness
states, not an explanation of why certain states are accompanied by
subjective experience.  A thermostat, a weather system, and a CPU
all have nonzero $\Sigma$ and $\tauasym$; they are presumably not
conscious.  The ``hard problem''~\cite{Chalmers1995} remains.

(b) That $\Phi$ (IIT) equals $\Sigma$ in general.  As discussed
in Sec.~\ref{subsec:IIT}, $\Phi$ measures spatial integration while
$\Sigma$ measures temporal irreversibility.  These are distinct
quantities with explicit counterexamples in both directions:
a reversible but integrated system has $\Phi > 0$, $\Sigma = 0$;
an irreversible but decomposable system has $\Sigma > 0$,
$\Phi = 0$.  The identification holds only in the quantum
extension (QIIT), and even there the precise equivalence
requires further work.

(c) The intermediate-$\Sigma$ and dream-state conjectures
(Secs.~\ref{sec:intermediate}--\ref{sec:dreams}).  These are
interpretive proposals motivated by the mathematical framework but
not derived from it.

%----------------------------------------------------------------------
\subsection{Relation to existing work}
\label{subsec:relation}
%----------------------------------------------------------------------

The idea that consciousness requires temporal irreversibility is not
new.  de~la~Fuente et al.~\cite{delaFuente2022} explicitly proposed
temporal irreversibility as a ``signature of consciousness,'' and
the FDT-violation framework~\cite{PRR2025} reaches similar
conclusions from the perspective of non-equilibrium statistical
mechanics.

What the present work adds is:

(i) A precise mathematical identification of the divergence measures
used across five independent theories (Table~\ref{tab:unification}).

(ii) A universal bound
($\tauasym \leq 1 - \exp(-\Sigma/2)$)
that constrains all five frameworks simultaneously, derived from a
theorem rather than an empirical correlation.

(iii) The identification of the Petz recovery map with Bayesian
retrodiction in the brain, connecting Friston's prediction-error
minimization with the mathematical structure of quantum recovery
theory.

(iv) A concrete, measurable quantity---the retrodiction complexity
index---that can be computed from existing neural data.

%----------------------------------------------------------------------
\subsection{Mathematical gaps and limitations}
\label{subsec:gaps}
%----------------------------------------------------------------------

Several mathematical gaps should be noted honestly.

\emph{The Markov assumption.}  Brain dynamics is not Markov.  Neural
activity exhibits long-range temporal correlations, $1/f$ spectral
scaling, and memory effects from synaptic plasticity.  The bound
(\ref{eq:brain_bound}) holds for Markov-approximated dynamics.  For
non-Markovian processes, a Markov embedding on an augmented state
space preserves the bound's validity but may loosen it
significantly.  Extension of the JRSWW bound to non-Markovian
process tensors~\cite{Pollock2018} is an open problem.

\emph{Coarse-graining.}  Brain entropy production is measured at the
macroscopic scale (fMRI voxels with $\sim$3~mm resolution, or
$\sim$100 EEG channels).  By the data-processing inequality,
$\Sbrain^{\mathrm{coarse}} \leq \Sbrain^{\mathrm{micro}}$.  The
macroscopic bound is self-consistent at its own scale but
underestimates the true microscopic entropy production.

\emph{Classical vs.\ quantum.}  The paper treats brain dynamics as
classical, which is appropriate for the macroscopic scales of fMRI
and EEG.  If quantum coherence in microtubules is
relevant~\cite{Barbatti2024}---as Penrose and Hameroff
propose---the quantum version of the bound would apply and could be
tighter due to entanglement contributions.  Current evidence is
insufficient to resolve this question.

\emph{The hard problem.}  Nothing in this paper addresses why
certain physical processes are accompanied by subjective experience.
The framework identifies mathematical constraints on which physical
states \emph{can} be conscious (if consciousness requires temporal
irreversibility), not why they \emph{are}.

%----------------------------------------------------------------------
\subsection{Conclusion}
\label{subsec:conclusion}
%----------------------------------------------------------------------

Five theories of consciousness, developed independently over three
decades, share a common mathematical core: each places a
form of relative entropy at its center.  The Petz recovery bound
$\tauasym \leq 1 - \exp(-\Sigma/2)$ connects all five through a
single inequality, establishing a thermodynamic floor on the
quality of retrodiction---and hence, we argue, on the temporal
coherence of experience.

The brain, in this view, is a retrodiction machine.  It constructs
a model of its own past from its present state, and the fidelity of
this construction is bounded by the entropy it produces.
Consciousness exists in the regime where this entropy production
is large enough to sustain a meaningful time arrow yet structured
enough that retrodiction remains possible.  Below this regime lies
dreamless sleep; above it, perhaps, the dissolution of structured
experience.

Whether retrodiction fidelity is merely correlated with
consciousness or constitutive of it remains an open question.  But
the mathematical unification itself---that Friston, Tononi,
Sanz~Perl, Penrose, and Carhart-Harris are all measuring the same
divergence---is a theorem, not a speculation.  It is a starting
point for the quantitative science of temporal experience.


%======================================================================
% Acknowledgments
%======================================================================
\begin{acknowledgments}
The author thanks the contributors to the open quantum information
theory literature whose work made this synthesis possible.
\end{acknowledgments}


%======================================================================
% Bibliography
%======================================================================
\begin{thebibliography}{40}

\bibitem{Friston2019}
K.~Friston,
``A free energy principle for a particular physics,''
arXiv:1906.10184 (2019).

\bibitem{Albantakis2023}
L.~Albantakis, L.~Barbosa, G.~Findlay, et al.,
``Integrated information theory (IIT) 4.0: Formulating the properties
of phenomenal existence in physical terms,''
PLoS Comput.\ Biol.\ \textbf{19}, e1011465 (2023).

\bibitem{Zanardi2018}
P.~Zanardi, M.~Tomka, and L.~C.~Venuti,
``Towards quantum integrated information theory,''
arXiv:1806.01421 (2018).

\bibitem{Albantakis2023Q}
L.~Albantakis, J.~Prentner, and I.~Durham,
``Computing the integrated information of a quantum mechanism,''
Entropy \textbf{25}, 449 (2023).

\bibitem{SanzPerl2021}
Y.~Sanz~Perl, C.~Pallavicini, I.~P\'{e}rez~Ipi\~{n}a, et al.,
``Nonequilibrium brain dynamics as a signature of consciousness,''
Phys.\ Rev.\ E \textbf{104}, 014411 (2021).

\bibitem{Penrose1996}
R.~Penrose,
``On gravity's role in quantum state reduction,''
Gen.\ Relativ.\ Gravit.\ \textbf{28}, 581 (1996).

\bibitem{Barbatti2024}
M.~Barbatti, L.~Gonz\'{a}lez, and M.~Persico,
``Gravitational decoherence and collapse time for molecules,''
J.\ Chem.\ Phys.\ (2024).

\bibitem{CarhartHarris2014}
R.~L.~Carhart-Harris, R.~Leech, P.~J.~Hellyer, et al.,
``The entropic brain: A theory of conscious states informed by
neuroimaging research with psychedelic drugs,''
Front.\ Hum.\ Neurosci.\ \textbf{8}, 20 (2014).

\bibitem{CarhartHarris2018}
R.~L.~Carhart-Harris,
``The entropic brain---revisited,''
Neuropharmacology \textbf{142}, 167 (2018).

\bibitem{Petz1986}
D.~Petz,
``Sufficient subalgebras and the relative entropy of states of a
von~Neumann algebra,''
Commun.\ Math.\ Phys.\ \textbf{105}, 123 (1986).

\bibitem{Petz1988}
D.~Petz,
``Sufficiency of channels over von~Neumann algebras,''
Q.\ J.\ Math.\ \textbf{39}, 97 (1988).

\bibitem{Parzygnat2023}
A.~J.~Parzygnat and F.~Buscemi,
``Axioms for retrodiction: Achieving time-reversal symmetry with
a prior,''
Quantum \textbf{7}, 1013 (2023).

\bibitem{Junge2018}
M.~Junge, R.~Renner, D.~Sutter, M.~M.~Wilde, and A.~Winter,
``Universal recovery maps and approximate sufficiency of quantum
relative entropy,''
Ann.\ Henri Poincar\'{e} \textbf{19}, 2955 (2018).

\bibitem{HuangPaper1}
S.-K.~Huang,
``Temporal asymmetry from the Petz recovery map: Entropy production,
the arrow of time, and the chain rule of quantum relative entropy,''
(2026), \url{https://github.com/akaiHuang/petz-recovery-unification}.

\bibitem{HuangPaper1b}
S.-K.~Huang,
``Quantum collapse as Petz recovery failure: An information-theoretic
resolution,''
(2026).

\bibitem{Fields2022}
C.~Fields, K.~Friston, J.~F.~Glazebrook, and M.~Levin,
``A free energy principle for generic quantum systems,''
Prog.\ Biophys.\ Mol.\ Biol.\ \textbf{173}, 36 (2022).

\bibitem{Casali2013}
A.~G.~Casali, O.~Gosseries, M.~Rosanova, et al.,
``A theoretically based index of consciousness independent of
sensory processing and behavior,''
Sci.\ Transl.\ Med.\ \textbf{5}, 198ra105 (2013).

\bibitem{Schartner2017}
M.~M.~Schartner, A.~Pigorini, S.~A.~Gibbs, et al.,
``Global and local complexity of intracranial EEG decreases during
NREM sleep,''
Neurosci.\ Conscious.\ \textbf{2017}, niw022 (2017).

\bibitem{delaFuente2022}
L.~de~la~Fuente, F.~Zamberlan, H.~Bocaccio, et al.,
``Temporal irreversibility of neural dynamics as a signature of
consciousness,''
Cereb.\ Cortex \textbf{33}, 1856 (2023).

\bibitem{Lynn2023}
C.~W.~Lynn, E.~J.~Cornblath, L.~Papadopoulos, M.~A.~Bertolero,
and D.~S.~Bassett,
``Entropy production of multivariate Ornstein-Uhlenbeck processes
correlates with consciousness levels in the human brain,''
Phys.\ Rev.\ E \textbf{107}, 024121 (2023).

\bibitem{PRR2025}
Y.~Sanz~Perl, et al.,
``Thermodynamics of consciousness: Violations of the
fluctuation-dissipation theorem as markers of conscious states,''
Phys.\ Rev.\ Research \textbf{7}, 013301 (2025).

\bibitem{Schnakenberg1976}
J.~Schnakenberg,
``Network theory of microscopic and macroscopic behavior of master
equation systems,''
Rev.\ Mod.\ Phys.\ \textbf{48}, 571 (1976).

\bibitem{Voss2014}
U.~Voss, R.~Holzmann, A.~Hobson, et al.,
``Induction of self awareness in dreams through frontal low current
stimulation of gamma activity,''
Nat.\ Neurosci.\ \textbf{17}, 810 (2014).

\bibitem{Erlacher2014}
D.~Erlacher, T.~Sch\"{a}dlich, T.~Stumbrys, and M.~Schredl,
``Time for actions in lucid dreams: Effects of task modality, length,
and complexity,''
Front.\ Psychol.\ \textbf{4}, 1013 (2014).

\bibitem{NIH2019}
T.~Izawa, S.~Yamashita, et al.,
``REM sleep is made of hard-won wakefulness,''
Cell \textbf{181}, 1346 (2019).

\bibitem{Chalmers1995}
D.~J.~Chalmers,
``Facing up to the problem of consciousness,''
J.\ Conscious.\ Stud.\ \textbf{2}, 200 (1995).

\bibitem{Pollock2018}
F.~A.~Pollock, C.~Rodr\'{i}guez-Rosario, T.~Frauenheim,
M.~Paternostro, and K.~Modi,
``Non-Markovian quantum processes: Complete framework and efficient
characterization,''
Phys.\ Rev.\ A \textbf{97}, 012127 (2018).

\bibitem{Landi2021}
G.~T.~Landi and M.~Paternostro,
``Irreversible entropy production: From classical to quantum,''
Rev.\ Mod.\ Phys.\ \textbf{93}, 035008 (2021).

\bibitem{FawziRenner2015}
O.~Fawzi and R.~Renner,
``Quantum conditional mutual information and approximate Markov
chains,''
Commun.\ Math.\ Phys.\ \textbf{340}, 575 (2015).

\bibitem{Phillips2024}
S.~Phillips, et al.,
``Integration, information, and exclusion correspond to limits,
colimits, and adjunctions,''
(2024).

\bibitem{Kleiner2020}
J.~Kleiner and S.~Tull,
``The mathematical structure of integrated information theory,''
Front.\ Appl.\ Math.\ Stat.\ \textbf{6}, 602973 (2021).

\bibitem{Ramstead2023}
M.~J.~D.~Ramstead, D.~A.~R.~Sakthivadivel,
C.~Heins, et al.,
``On Bayesian mechanics: A physics of and by beliefs,''
Interface Focus \textbf{13}, 20220029 (2023).

\bibitem{Hameroff2025}
S.~Hameroff, A.~Khrennikov, et al.,
``Conscious active inference~II: Quantum orchestrated objective
reduction among intraneuronal microtubules naturally accounts for
discrete perceptual cycles,''
Biosystems (2025).

\bibitem{DoerigUnfolding}
A.~Doerig, A.~Schurger, K.~Hess, and M.~H.~Herzog,
``The unfolding argument: Why IIT and other causal structure theories
cannot explain consciousness,''
Conscious.\ Cogn.\ \textbf{72}, 49 (2019).

\bibitem{Buscemi2024}
F.~Buscemi, J.~Schindler, and D.~Sutter,
``Approximate recoverability and relative entropy,''
arXiv:2412.12489 (2024).

\end{thebibliography}

\end{document}
